\documentclass[10pt,a4paper]{amsart}
\usepackage{amsmath,amssymb,amsthm,graphicx}
% Retained minimal preamble for the Stacks Project chapter sites-cohomology.tex (Cohomology on Sites).
% Verbatim excerpt of the chapter preamble: Xy-pic setup (native preamble lines 14-18) and the native
% theorem declarations with the mathematical macros (native preamble lines 154-217). Copied mechanically;
% nothing added, removed or redefined. The wrapper class and the amsmath/amssymb/amsthm/graphicx base
% packages are supplied by the assembler, so the class/externaldocument part of the native preamble is omitted.
% For commutative diagrams we use Xy-pic
\usepackage[all]{xy}

% We use 2cell for 2-commutative diagrams.
\xyoption{2cell}
% Theorem environments.
%
\theoremstyle{plain}
\newtheorem{theorem}[subsection]{Theorem}
\newtheorem{proposition}[subsection]{Proposition}
\newtheorem{lemma}[subsection]{Lemma}

\theoremstyle{definition}
\newtheorem{definition}[subsection]{Definition}
\newtheorem{example}[subsection]{Example}
\newtheorem{exercise}[subsection]{Exercise}
\newtheorem{situation}[subsection]{Situation}

\theoremstyle{remark}
\newtheorem{remark}[subsection]{Remark}
\newtheorem{remarks}[subsection]{Remarks}

\numberwithin{equation}{subsection}

% Macros
%
\def\lim{\mathop{\mathrm{lim}}\nolimits}
\def\colim{\mathop{\mathrm{colim}}\nolimits}
\def\Spec{\mathop{\mathrm{Spec}}}
\def\Hom{\mathop{\mathrm{Hom}}\nolimits}
\def\Ext{\mathop{\mathrm{Ext}}\nolimits}
\def\SheafHom{\mathop{\mathcal{H}\!\mathit{om}}\nolimits}
\def\SheafExt{\mathop{\mathcal{E}\!\mathit{xt}}\nolimits}
\def\Sch{\mathit{Sch}}
\def\Mor{\mathop{\mathrm{Mor}}\nolimits}
\def\Ob{\mathop{\mathrm{Ob}}\nolimits}
\def\Sh{\mathop{\mathit{Sh}}\nolimits}
\def\NL{\mathop{N\!L}\nolimits}
\def\CH{\mathop{\mathrm{CH}}\nolimits}
\def\proetale{{pro\text{-}\acute{e}tale}}
\def\etale{{\acute{e}tale}}
\def\QCoh{\mathit{QCoh}}
\def\Ker{\mathop{\mathrm{Ker}}}
\def\Im{\mathop{\mathrm{Im}}}
\def\Coker{\mathop{\mathrm{Coker}}}
\def\Coim{\mathop{\mathrm{Coim}}}

% Boxtimes
%
\DeclareMathSymbol{\boxtimes}{\mathbin}{AMSa}{"02}

%
% Macros for moduli stacks/spaces
%
\def\QCohstack{\mathcal{QC}\!\mathit{oh}}
\def\Cohstack{\mathcal{C}\!\mathit{oh}}
\def\Spacesstack{\mathcal{S}\!\mathit{paces}}
\def\Quotfunctor{\mathrm{Quot}}
\def\Hilbfunctor{\mathrm{Hilb}}
\def\Curvesstack{\mathcal{C}\!\mathit{urves}}
\def\Polarizedstack{\mathcal{P}\!\mathit{olarized}}
\def\Complexesstack{\mathcal{C}\!\mathit{omplexes}}
% \Pic is the operator that assigns to X its picard group, usage \Pic(X)
% \Picardstack_{X/B} denotes the Picard stack of X over B
% \Picardfunctor_{X/B} denotes the Picard functor of X over B
\def\Pic{\mathop{\mathrm{Pic}}\nolimits}
\def\Picardstack{\mathcal{P}\!\mathit{ic}}
\def\Picardfunctor{\mathrm{Pic}}
\def\Deformationcategory{\mathcal{D}\!\mathit{ef}}

% --- appended 2026-10-12 (Chief): the project own definitions of the reference/slogan/history
% display environments, copied mechanically from the shared Stacks preamble lines 8-12
% (\usepackage{verbatim} plus the three \newenvironment lines, which the chapter preamble keeps
% verbatim from the project preamble). Nothing added, removed or redefined: these environments
% carry no mathematical content and the project itself typesets them as comments. Without these
% three lines a counted statement whose own text uses them fails with Environment slogan undefined.
\usepackage{verbatim}
\newenvironment{reference}{\comment}{\endcomment}
\newenvironment{slogan}{\comment}{\endcomment}
\newenvironment{history}{\comment}{\endcomment}
\usepackage{cleveref}
\crefname{originalequation}{original equation}{original equations}
\Crefname{originalequation}{Original equation}{Original equations}
\makeatletter
\newlabel{equation-internal-hom}{{08DI}{source}{}{sourcepointer-1}{}}
\newlabel{equation-internal-hom@cref}{{[originalequation][0][]08DI}{[1][0][]source}{}{}{}}
\makeatother
\crefname{originalreference}{original reference}{original references}
\Crefname{originalreference}{Original reference}{Original references}
\makeatletter
\newlabel{lemma-RHom-from-K-flat-into-K-injective}{{0A8T}{source}{}{sourcepointer-2}{}}
\newlabel{lemma-RHom-from-K-flat-into-K-injective@cref}{{[originalreference][0][]0A8T}{[1][0][]source}{}{}{}}
\makeatother
\crefname{originalreference}{original reference}{original references}
\Crefname{originalreference}{Original reference}{Original references}
\makeatletter
\newlabel{lemma-compose}{{0A5Y}{source}{}{sourcepointer-3}{}}
\newlabel{lemma-compose@cref}{{[originalreference][0][]0A5Y}{[1][0][]source}{}{}{}}
\makeatother
\crefname{originalsection}{original section}{original sections}
\Crefname{originalsection}{Original section}{Original sections}
\makeatletter
\newlabel{lemma-section-RHom-over-U}{{08DK}{source}{}{sourcepointer-4}{}}
\newlabel{lemma-section-RHom-over-U@cref}{{[originalsection][0][]08DK}{[1][0][]source}{}{}{}}
\makeatother
\title{The internal hom into an internal hom is the internal hom from the derived tensor product}
\author{The Stacks Project Authors}
\date{stacks-project:sites-cohomology.tex:local-snapshot-20261011; GNU FDL 1.2 or later}
\begin{document}
\maketitle
\noindent\textit{Editorial note.} All delivered mathematical text is the author's own, copied from the pinned Stacks Project chapter file sites-cohomology.tex; no statement or proof step was rewritten, reordered or removed.\par\medskip
\noindent\textit{Editorial note.} The counted proof is the author's complete argument: the proof is written out in full: the choice of representatives, the displayed identity of complexes from the composition lemma, the identification of the two sides as derived internal homs, and the recovery of the chapter's global-sections equation.\par\medskip
\noindent\textit{Editorial note.} The label lemma-internal-hom names the counted statement, and the proof environment immediately following it in the source is delivered verbatim as its proof.\par\medskip
\noindent\textit{Editorial note.} The counted statement is carried by the author's own lemma environment in the chapter source.\par\medskip
\noindent\textit{Editorial note.} This card belongs to the priority queue of the D-72 twin ruling (the sites-side statement is the designated counting side, now counted here for the first time). The statement and proof are copied from the pinned chapter file; the three cited inputs are kept as external pointers with their canonical Stacks tags.\par\medskip
\noindent\textit{Allowed standard prerequisites (editorial).}\begin{itemize}
\item Composition of internal homs. Same chapter, the lemma identifying the complex of internal homs into an internal hom with the internal hom from the total tensor product. Supplies the displayed identity of complexes at the heart of the proof.
\item The internal hom from a K-flat complex into a K-injective complex is K-injective. Same chapter, the lemma stating that the internal hom from a K-flat complex into a K-injective complex is K-injective. Supplies that the left-hand complex represents the derived internal hom into the derived internal hom.
\item Zeroth cohomology of the derived internal hom as derived Hom. Same chapter, the lemma identifying H\textasciicircum{}0 of the internal hom complex with the derived Hom groups over an object and globally. Supplies the identification of the global sections formula after taking H\textasciicircum{}0.
\end{itemize}\medskip
\ifdefined\hypertarget\hypertarget{sourcepointer-1}{}\fi
\noindent\textit{Source reference.} Original equation 08DI: sources/stacks-sites-cohomology/pinned-original-source.tex, source lines 8844--8844. Prerequisite pointer invoked at proof line(s) 8905; it names a standard fact of the same chapter that the counted argument uses, and every step of the counted proof built on it is retained verbatim in the delivered body.\par\medskip
\ifdefined\hypertarget\hypertarget{sourcepointer-2}{}\fi
\noindent\textit{Source reference.} Original reference 0A8T: sources/stacks-sites-cohomology/pinned-original-source.tex, source lines 8789--8789. Prerequisite pointer invoked at proof line(s) 8925; it names a standard fact of the same chapter that the counted argument uses, and every step of the counted proof built on it is retained verbatim in the delivered body.\par\medskip
\ifdefined\hypertarget\hypertarget{sourcepointer-3}{}\fi
\noindent\textit{Source reference.} Original reference 0A5Y: sources/stacks-sites-cohomology/pinned-original-source.tex, source lines 8616--8616. Prerequisite pointer invoked at proof line(s) 8922; it names a standard fact of the same chapter that the counted argument uses, and every step of the counted proof built on it is retained verbatim in the delivered body.\par\medskip
\ifdefined\hypertarget\hypertarget{sourcepointer-4}{}\fi
\noindent\textit{Source reference.} Original section 08DK: sources/stacks-sites-cohomology/pinned-original-source.tex, source lines 8874--8874. Prerequisite pointer invoked at proof line(s) 8929; it names a standard fact of the same chapter that the counted argument uses, and every step of the counted proof built on it is retained verbatim in the delivered body.\par\medskip
\begin{lemma}
\label{lemma-internal-hom}
Let $(\mathcal{C}, \mathcal{O})$ be a ringed site. Let $K, L, M$ be objects
of $D(\mathcal{O})$. With the construction as described above
there is a canonical isomorphism
$$
R\SheafHom(K, R\SheafHom(L, M)) =
R\SheafHom(K \otimes_\mathcal{O}^\mathbf{L} L, M)
$$
in $D(\mathcal{O})$ functorial in $K, L, M$
which recovers (\ref{equation-internal-hom}) on taking $H^0(\mathcal{C}, -)$.
\end{lemma}

\begin{proof}
Choose a K-injective complex $\mathcal{I}^\bullet$ representing
$M$ and a K-flat complex of $\mathcal{O}$-modules $\mathcal{L}^\bullet$
representing $L$.
For any complex of $\mathcal{O}$-modules $\mathcal{K}^\bullet$
we have
$$
\SheafHom^\bullet(\mathcal{K}^\bullet,
\SheafHom^\bullet(\mathcal{L}^\bullet, \mathcal{I}^\bullet))
=
\SheafHom^\bullet(
\text{Tot}(\mathcal{K}^\bullet \otimes_\mathcal{O} \mathcal{L}^\bullet),
\mathcal{I}^\bullet)
$$
by Lemma \ref{lemma-compose}.
Note that the left hand side represents
$R\SheafHom(K, R\SheafHom(L, M))$ (use
Lemma \ref{lemma-RHom-from-K-flat-into-K-injective})
and that the right hand side represents
$R\SheafHom(K \otimes_\mathcal{O}^\mathbf{L} L, M)$.
This proves the displayed formula of the lemma.
Taking global sections and using Lemma \ref{lemma-section-RHom-over-U}
we obtain (\ref{equation-internal-hom}).
\end{proof}

\end{document}
